\section{Sliding Window Distinct Values}
\label{sec:cses-sw-distinct}

\problemstatement{Given an array of $n$ integers and a window size $k$,
output the number of distinct values in every contiguous window of size
$k$.}

\begin{patternbox}
Distinct-count slides cleanly with a \textbf{frequency hash map}: track how
many times each value appears in the window. A value contributes to the
distinct count exactly while its frequency is positive, so the count
changes only when a frequency crosses between $0$ and $1$.
\end{patternbox}

\subsection*{Frequency map with a distinct counter}

Keep $\textit{freq}[v]$ and a running $\textit{distinct}$. Adding a value
$v$: if its frequency was $0$, it is newly present, so increment
$\textit{distinct}$; then increment $\textit{freq}[v]$. Removing a value
$v$: decrement $\textit{freq}[v]$; if it drops to $0$, the value has left
the window, so decrement $\textit{distinct}$. Each slide performs one add
and one remove, both $O(1)$ amortised.

\begin{ideabox}[Watch the Zero Boundary]
The distinct count only moves when a frequency transitions $0\to1$ (a value
appears) or $1\to0$ (a value disappears). Checking that boundary on each
add/remove keeps the counter exact without rescanning the window.
\end{ideabox}

\subsection*{Algorithm}

\begin{enumerate}
  \item Add the first $k$ elements, updating $\textit{freq}$ and
        $\textit{distinct}$.
  \item For each slide: add $a[i]$ (increment distinct if it was absent),
        remove $a[i-k]$ (decrement distinct if it becomes absent).
  \item Output $\textit{distinct}$ per window.
\end{enumerate}

\subsection*{Dry run}

\dryrun{$a=[1,2,1,3]$, $k=3$.}

\begin{itemize}
  \item Window $[1,2,1]$: freqs $\{1{:}2,2{:}1\}$, distinct $=2$.
  \item Slide to $[2,1,3]$: remove a $1$ (freq $2{\to}1$, still present),
        add $3$ (new) $\Rightarrow$ distinct $=3$.
  \item Output $2,3$.
\end{itemize}

\subsection*{C++ implementation}

\begin{lstlisting}
#include <bits/stdc++.h>
using namespace std;

int main() {
    int n, k;
    cin >> n >> k;
    vector<int> a(n);
    for (int& x : a) cin >> x;

    unordered_map<int,int> freq;
    int distinct = 0;
    auto add = [&](int v) { if (freq[v]++ == 0) distinct++; };
    auto remove = [&](int v) { if (--freq[v] == 0) distinct--; };

    for (int i = 0; i < k; i++) add(a[i]);
    cout << distinct << " ";
    for (int i = k; i < n; i++) {
        add(a[i]); remove(a[i - k]);
        cout << distinct << " ";
    }
    cout << "\n";
    return 0;
}
\end{lstlisting}

\begin{complexitybox}
\textbf{Time Complexity.} $O(n)$ amortised with a hash map.
\textbf{Space Complexity.} $O(k)$ distinct keys at most.
\end{complexitybox}

\begin{takeawaybox}
Distinct-value counting is a frequency map plus a counter that only reacts
at the zero boundary. This add/remove-with-boundary-check idiom generalises
to ``count of values with property P'' whenever P depends on a per-value
frequency.
\end{takeawaybox}
